\section{Num\-Mat$<$ F $>$ Class Template Reference}
\label{classNumMat}\index{NumMat@{NumMat}}
{\tt \#include $<$nummat.hpp$>$}



\subsection{Detailed Description}
\subsubsection*{template$<$class F$>$ class Num\-Mat$<$ F $>$}

Num\-Mat is a matrix made up of a size and data. The data can be made up of any type (ints, double, bools, etc.). 



Definition at line 27 of file nummat.hpp.\subsection*{Public Member Functions}
\begin{CompactItemize}
\item 
{\bf Num\-Mat} (int m=0, int n=0)
\item 
{\bf Num\-Mat} (int m, int n, bool owndata, F $\ast$data)
\item 
{\bf Num\-Mat} (const {\bf Num\-Mat} \&C)
\item 
{\bf $\sim$Num\-Mat} ()
\item 
{\bf Num\-Mat} \& {\bf operator=} (const {\bf Num\-Mat} \&C)
\item 
void {\bf resize} (int m, int n)
\item 
const F \& {\bf operator()} (int i, int j) const 
\item 
F \& {\bf operator()} (int i, int j)
\item 
F $\ast$ {\bf data} () const 
\item 
F $\ast$ {\bf clmdata} (int j)
\item 
int {\bf m} () const 
\item 
int {\bf n} () const 
\end{CompactItemize}
\subsection*{Public Attributes}
\begin{CompactItemize}
\item 
int {\bf \_\-m}
\item 
int {\bf \_\-n}
\item 
bool {\bf \_\-owndata}
\item 
F $\ast$ {\bf \_\-data}
\end{CompactItemize}


\subsection{Constructor \& Destructor Documentation}
\index{NumMat@{Num\-Mat}!NumMat@{NumMat}}
\index{NumMat@{NumMat}!NumMat@{Num\-Mat}}
\subsubsection{\setlength{\rightskip}{0pt plus 5cm}template$<$class F$>$ {\bf Num\-Mat}$<$ F $>$::{\bf Num\-Mat} (int {\em m} = {\tt 0}, int {\em n} = {\tt 0})\hspace{0.3cm}{\tt  [inline]}}\label{classNumMat_dbf73e11c1ff70b095c6d97aa75ed610_dbf73e11c1ff70b095c6d97aa75ed610}


Construct a Num\-Mat using the dimensions, m and n. owndata set to true 

Definition at line 40 of file nummat.hpp.

References Num\-Mat$<$ F $>$::\_\-data, Num\-Mat$<$ F $>$::\_\-m, and Num\-Mat$<$ F $>$::\_\-n.

\footnotesize\begin{verbatim}40                           : _m(m), _n(n), _owndata(true) {
41          if(_m>0 && _n>0) { _data = new F[_m*_n]; assert( _data!=NULL ); memset(_data, 0, _m*_n*sizeof(F)); } else _data=NULL;
42   }
\end{verbatim}\normalsize 


\index{NumMat@{Num\-Mat}!NumMat@{NumMat}}
\index{NumMat@{NumMat}!NumMat@{Num\-Mat}}
\subsubsection{\setlength{\rightskip}{0pt plus 5cm}template$<$class F$>$ {\bf Num\-Mat}$<$ F $>$::{\bf Num\-Mat} (int {\em m}, int {\em n}, bool {\em owndata}, F $\ast$ {\em data})\hspace{0.3cm}{\tt  [inline]}}\label{classNumMat_b5da1f826a1fdb2179102482afe8981d_b5da1f826a1fdb2179102482afe8981d}


Construct Num\-Mat with data given. If owndata is false, then the data passed is not stored in a new array, F. Otherwise, an array of type F must be allocated into memory to store the data 

Definition at line 45 of file nummat.hpp.

\footnotesize\begin{verbatim}45                                              : _m(m), _n(n), _owndata(owndata) {
46          if(_owndata) {
47                 if(_m>0 && _n>0) { _data = new F[_m*_n]; assert( _data!=NULL ); memset(_data, 0, _m*_n*sizeof(F)); } else _data=NULL;
48                 if(_m>0 && _n>0) memcpy( _data, data, _m*_n*sizeof(F) );
49          } else {
50                 _data = data;
51          }
52   }
\end{verbatim}\normalsize 


\index{NumMat@{Num\-Mat}!NumMat@{NumMat}}
\index{NumMat@{NumMat}!NumMat@{Num\-Mat}}
\subsubsection{\setlength{\rightskip}{0pt plus 5cm}template$<$class F$>$ {\bf Num\-Mat}$<$ F $>$::{\bf Num\-Mat} (const {\bf Num\-Mat}$<$ F $>$ \& {\em C})\hspace{0.3cm}{\tt  [inline]}}\label{classNumMat_4b3436a830e943717b773e03eebf08b7_4b3436a830e943717b773e03eebf08b7}


Construct a new matrix with another matrix as argument. New matrix gets C.\_\-m, C.\_\-n and C.\_\-owndata. If \_\-owndata is true, the new matrix generates an array F to copy the data. If \_\-owndata is false, \_\-data is just set to C.\_\-data, simply pointing to the same data 

Definition at line 56 of file nummat.hpp.

References Num\-Mat$<$ F $>$::\_\-data.

\footnotesize\begin{verbatim}56                          : _m(C._m), _n(C._n), _owndata(C._owndata) {
57          if(_owndata) {
58                 if(_m>0 && _n>0) { _data = new F[_m*_n]; assert( _data!=NULL ); memset(_data, 0, _m*_n*sizeof(F)); } else _data=NULL;
59                 if(_m>0 && _n>0) memcpy( _data, C._data, _m*_n*sizeof(F) );
60          } else {
61                 _data = C._data;
62          }
63   }
\end{verbatim}\normalsize 


\index{NumMat@{Num\-Mat}!~NumMat@{$\sim$NumMat}}
\index{~NumMat@{$\sim$NumMat}!NumMat@{Num\-Mat}}
\subsubsection{\setlength{\rightskip}{0pt plus 5cm}template$<$class F$>$ {\bf Num\-Mat}$<$ F $>$::$\sim${\bf Num\-Mat} ()\hspace{0.3cm}{\tt  [inline]}}\label{classNumMat_2de3b0b2c141a45e3738519e8732710a_2de3b0b2c141a45e3738519e8732710a}


Destroy Num\-Mat. If \_\-owndata is true, \_\-data must be deleted and deallocated 

Definition at line 65 of file nummat.hpp.

\footnotesize\begin{verbatim}65             { 
66          if(_owndata) { 
67                 if(_m>0 && _n>0) { delete[] _data; _data = NULL; } 
68          }
69   }
\end{verbatim}\normalsize 




\subsection{Member Function Documentation}
\index{NumMat@{Num\-Mat}!clmdata@{clmdata}}
\index{clmdata@{clmdata}!NumMat@{Num\-Mat}}
\subsubsection{\setlength{\rightskip}{0pt plus 5cm}template$<$class F$>$ F$\ast$ {\bf Num\-Mat}$<$ F $>$::clmdata (int {\em j})\hspace{0.3cm}{\tt  [inline]}}\label{classNumMat_29a94d113c6ed3f2191bb35c48a46e16_29a94d113c6ed3f2191bb35c48a46e16}


Return \_\-data j$\ast$\_\-m 

Definition at line 106 of file nummat.hpp.

\footnotesize\begin{verbatim}106 { return &(_data[j*_m]); }
\end{verbatim}\normalsize 


\index{NumMat@{Num\-Mat}!data@{data}}
\index{data@{data}!NumMat@{Num\-Mat}}
\subsubsection{\setlength{\rightskip}{0pt plus 5cm}template$<$class F$>$ F$\ast$ {\bf Num\-Mat}$<$ F $>$::data () const\hspace{0.3cm}{\tt  [inline]}}\label{classNumMat_ad3b2b9598c8547b4f7afe7a79ed520c_ad3b2b9598c8547b4f7afe7a79ed520c}


Return \_\-data 

Definition at line 104 of file nummat.hpp.

Referenced by clear(), SVDRep::construct(), daxpy(), dgemm(), dgemv(), SVDRep::dgemv(), dger(), and inv().

\footnotesize\begin{verbatim}104 { return _data; }
\end{verbatim}\normalsize 


\index{NumMat@{Num\-Mat}!m@{m}}
\index{m@{m}!NumMat@{Num\-Mat}}
\subsubsection{\setlength{\rightskip}{0pt plus 5cm}template$<$class F$>$ int {\bf Num\-Mat}$<$ F $>$::m () const\hspace{0.3cm}{\tt  [inline]}}\label{classNumMat_6a2fa69c25bf8ed36e019d0fb3f5e966_6a2fa69c25bf8ed36e019d0fb3f5e966}


Return m 

Definition at line 108 of file nummat.hpp.

Referenced by clear(), SVDRep::construct(), daxpy(), dgemm(), dgemv(), SVDRep::dgemv(), dger(), inv(), SVDRep::m(), operator$<$$<$(), pinv(), setvalue(), and tran().

\footnotesize\begin{verbatim}108 { return _m; }
\end{verbatim}\normalsize 


\index{NumMat@{Num\-Mat}!n@{n}}
\index{n@{n}!NumMat@{Num\-Mat}}
\subsubsection{\setlength{\rightskip}{0pt plus 5cm}template$<$class F$>$ int {\bf Num\-Mat}$<$ F $>$::n () const\hspace{0.3cm}{\tt  [inline]}}\label{classNumMat_611a0b095fab035a528671dd8cf7f3b7_611a0b095fab035a528671dd8cf7f3b7}


Return n 

Definition at line 110 of file nummat.hpp.

Referenced by clear(), SVDRep::construct(), daxpy(), dgemm(), dgemv(), SVDRep::dgemv(), dger(), inv(), SVDRep::n(), operator$<$$<$(), pinv(), setvalue(), and tran().

\footnotesize\begin{verbatim}110 { return _n; }
\end{verbatim}\normalsize 


\index{NumMat@{Num\-Mat}!operator()@{operator()}}
\index{operator()@{operator()}!NumMat@{Num\-Mat}}
\subsubsection{\setlength{\rightskip}{0pt plus 5cm}template$<$class F$>$ F\& {\bf Num\-Mat}$<$ F $>$::operator() (int {\em i}, int {\em j})\hspace{0.3cm}{\tt  [inline]}}\label{classNumMat_743eed4071652374b372885106b9142f_743eed4071652374b372885106b9142f}


Return \_\-data(i,j) or \_\-data[i+j$\ast$\_\-m] 

Definition at line 99 of file nummat.hpp.

\footnotesize\begin{verbatim}99                                { 
100          assert( i>=0 && i<_m && j>=0 && j<_n );
101          return _data[i+j*_m];
102   }
\end{verbatim}\normalsize 


\index{NumMat@{Num\-Mat}!operator()@{operator()}}
\index{operator()@{operator()}!NumMat@{Num\-Mat}}
\subsubsection{\setlength{\rightskip}{0pt plus 5cm}template$<$class F$>$ const F\& {\bf Num\-Mat}$<$ F $>$::operator() (int {\em i}, int {\em j}) const\hspace{0.3cm}{\tt  [inline]}}\label{classNumMat_08092f91032a8a0df0ad0870710e6a7a_08092f91032a8a0df0ad0870710e6a7a}


Return \_\-data(i,j) or \_\-data[i+j$\ast$\_\-m] 

Definition at line 94 of file nummat.hpp.

\footnotesize\begin{verbatim}94                                            { 
95          assert( i>=0 && i<_m && j>=0 && j<_n );
96          return _data[i+j*_m];
97   }
\end{verbatim}\normalsize 


\index{NumMat@{Num\-Mat}!operator=@{operator=}}
\index{operator=@{operator=}!NumMat@{Num\-Mat}}
\subsubsection{\setlength{\rightskip}{0pt plus 5cm}template$<$class F$>$ {\bf Num\-Mat}\& {\bf Num\-Mat}$<$ F $>$::operator= (const {\bf Num\-Mat}$<$ F $>$ \& {\em C})\hspace{0.3cm}{\tt  [inline]}}\label{classNumMat_7e5baac1fa9ea6a2167d746fb6c83a4c_7e5baac1fa9ea6a2167d746fb6c83a4c}


A way of setting an existing matrix equal to another exisiting matrix 

Definition at line 71 of file nummat.hpp.

\footnotesize\begin{verbatim}71                                      {
72          if(_owndata) { 
73                 if(_m>0 && _n>0) { delete[] _data; _data = NULL; } 
74          }
75          _m = C._m; _n=C._n; _owndata=C._owndata;
76          if(_owndata) {
77                 if(_m>0 && _n>0) { _data = new F[_m*_n]; assert( _data!=NULL ); memset(_data, 0, _m*_n*sizeof(F)); } else _data=NULL;
78                 if(_m>0 && _n>0) memcpy( _data, C._data, _m*_n*sizeof(F) );
79          } else {
80                 _data = C._data;
81          }
82          return *this;
83   }
\end{verbatim}\normalsize 


\index{NumMat@{Num\-Mat}!resize@{resize}}
\index{resize@{resize}!NumMat@{Num\-Mat}}
\subsubsection{\setlength{\rightskip}{0pt plus 5cm}template$<$class F$>$ void {\bf Num\-Mat}$<$ F $>$::resize (int {\em m}, int {\em n})\hspace{0.3cm}{\tt  [inline]}}\label{classNumMat_2ebfd442ba48ad3dfe61b1c0fc370599_2ebfd442ba48ad3dfe61b1c0fc370599}


Resize the matrix with new m and n. Asserts \_\-owndata is true so that new array of type F can be allocated into memory 

Definition at line 85 of file nummat.hpp.

Referenced by SVDRep::construct().

\footnotesize\begin{verbatim}85                              {
86          assert( _owndata==true );
87          if(_m!=m || _n!=n) {
88                 if(_m>0 && _n>0) { delete[] _data; _data = NULL; } 
89                 _m = m; _n = n;
90                 if(_m>0 && _n>0) { _data = new F[_m*_n]; assert( _data!=NULL ); memset(_data, 0, _m*_n*sizeof(F)); } else _data=NULL;
91          }
92   }
\end{verbatim}\normalsize 




\subsection{Member Data Documentation}
\index{NumMat@{Num\-Mat}!_data@{\_\-data}}
\index{_data@{\_\-data}!NumMat@{Num\-Mat}}
\subsubsection{\setlength{\rightskip}{0pt plus 5cm}template$<$class F$>$ F$\ast$ {\bf Num\-Mat}$<$ F $>$::{\bf \_\-data}}\label{classNumMat_7125153abbcb2e949a1b4f81cbe22af5_7125153abbcb2e949a1b4f81cbe22af5}


The data. Since using templates, \_\-data can be of any type 

Definition at line 37 of file nummat.hpp.

Referenced by Num\-Mat$<$ F $>$::Num\-Mat().\index{NumMat@{Num\-Mat}!_m@{\_\-m}}
\index{_m@{\_\-m}!NumMat@{Num\-Mat}}
\subsubsection{\setlength{\rightskip}{0pt plus 5cm}template$<$class F$>$ int {\bf Num\-Mat}$<$ F $>$::{\bf \_\-m}}\label{classNumMat_b781ae81625400d76d94420d3af2c230_b781ae81625400d76d94420d3af2c230}


\_\-m is the \char`\"{}m\char`\"{} dimension, or number of rows in the matrix 

Definition at line 31 of file nummat.hpp.

Referenced by Num\-Mat$<$ F $>$::Num\-Mat().\index{NumMat@{Num\-Mat}!_n@{\_\-n}}
\index{_n@{\_\-n}!NumMat@{Num\-Mat}}
\subsubsection{\setlength{\rightskip}{0pt plus 5cm}template$<$class F$>$ int {\bf Num\-Mat}$<$ F $>$::{\bf \_\-n}}\label{classNumMat_b628b96b5cc160f67acf6b4673cdd8bb_b628b96b5cc160f67acf6b4673cdd8bb}


\_\-n is the \char`\"{}n\char`\"{} dimension, or number of columns in the matrix 

Definition at line 33 of file nummat.hpp.

Referenced by Num\-Mat$<$ F $>$::Num\-Mat().\index{NumMat@{Num\-Mat}!_owndata@{\_\-owndata}}
\index{_owndata@{\_\-owndata}!NumMat@{Num\-Mat}}
\subsubsection{\setlength{\rightskip}{0pt plus 5cm}template$<$class F$>$ bool {\bf Num\-Mat}$<$ F $>$::{\bf \_\-owndata}}\label{classNumMat_4a613b2956a4a0ffff962e497a5d56e0_4a613b2956a4a0ffff962e497a5d56e0}


Whether or not data is own 

Definition at line 35 of file nummat.hpp.

The documentation for this class was generated from the following file:\begin{CompactItemize}
\item 
common/{\bf nummat.hpp}\end{CompactItemize}
